\printedpage{145}
% head: A. Grothendieck : Dérivateurs, Ch. XII, Édition des Universités de Montpellier II et Paris VII
% foot: Version 30 mars 2003

\msp{page 228}
So the question arises of the dual factorization of an arrow $f : X \longrightarrow Y$ in Cat into
\[
\begin{tikzcd}[column sep=huge]
X \arrow[r, hook, "W\text{-cofibration}", "i_f"'] & \tilde{Y} \arrow[r, "\in W", "p_f"'] & Y.
\end{tikzcd}
\]

The results of \textsc{Thomason} ensure that this exists, with $i$ an open Dwyer immersion, and $p$ a $W$-fibration - at least in the case $W = W_\infty$. I am going to try to see it directly, by dually copying the construction of \textsc{Cartan-Serre}.

So it is first a matter of factoring
\[
X \amalg X \xrightarrow{\ \mathrm{codiag}_X\ } X,
\]
and for this I choose a fixed interval
\[
I \ = \ \Phi_n \ = \ (0 \longrightarrow 1 \longleftarrow 2 \longrightarrow \cdots \longrightarrow 2n-1 \longleftarrow 2n)
\]
with
\[
n \geq 2
\]
so that the inclusion
\[
\varepsilon = \text{discrete category } \{0,1\} \xrightarrow{\partial} I, \qquad \partial(0) = 0, \quad \partial(1) = 2n,
\]
is an open Dwyer immersion. This property is stable under Cat-fibrant base change, in particular on taking a product
\[
\underbrace{X \times \varepsilon}_{\simeq X \coprod X} \xrightarrow{i_X} X \times I,
\]
% continuation of a footnote begun on the preceding page of the edition (no footnote mark on this page)
\begingroup\renewcommand{\thefootnote}{}\footnotetext{%
\[
\begin{tikzcd}[ampersand replacement=\&]
U \arrow[r, hook, "i"] \arrow[d, "p"'] \& X \arrow[d, "q"] \\
U' \arrow[r, hook, "i'"] \& X'
\end{tikzcd}
\]
with $i$ an open Dwyer immersion, then
\[
\begin{aligned}
i \in W \quad &\Longleftrightarrow \quad i' \in W \\
p \in W \quad &\Longleftrightarrow \quad q \in W,
\end{aligned}
\]
so I need W(7 bis) \textit{and} even W(7) for the second implication.
\begin{enumerate}
\item[c)] Open Dwyer immersions \texttt{[are]} stable under cobase change and under composition, contain the isomorphisms.
\end{enumerate}}\endgroup
\printedpage{146}
% head: A. Grothendieck : Dérivateurs, Ch. XII, Édition des Universités de Montpellier II et Paris VII
% foot: Version 30 mars 2003
\msp{page 229}
on the other hand
\[
X \times I \xrightarrow{\mathrm{pr}_1} X
\]
is $W_\omega$-perfect, whence the desired factorization % [sic] factoriastion -> factorisation, misprint
\[
X \amalg X \xrightarrow{i_X} X \times I \ \xrightarrow{\ p_X = \mathrm{pr}_1^X\ } X
\]
- nothing could be simpler! How to use it to construct the factorization of an arbitrary arrow $f : X \longrightarrow Y$, by dualizing the constructions relating to path spaces? We define $Z(f)$ by the cocartesian square in Cat
\[
\begin{tikzcd}[column sep=5em, row sep=4em]
X \arrow[r, dotted, "\mathrm{inj}_1^X"'] & X \times I \arrow[r, "f'"] & Z(f) \arrow[d, dotted, bend left=45, "\,p_f"] \\
& X \arrow[u, hook, "\mathrm{inj}_0^X"'] \arrow[r, "f"] & Y, \arrow[u, hook, "W"']
\end{tikzcd}
\]
we have
\[
Z(f) \xrightarrow{\alpha} Y \qquad \text{induced by } \mathrm{pr}_1 : X \times I \longrightarrow X,
\]
and the desired factorization of $f$ is
\[
X \ \xrightarrow{\ i_f = f' \circ \mathrm{inj}_1^X\ } Z(f) \xrightarrow{p_f} Y.
\]

I say that
\[
\left\{
\begin{aligned}
& i_f \text{ is an open Dwyer immersion, hence a } W\text{-cofibration.} \\
& p_f \in W.
\end{aligned}
\right.
\]

\msp{page 230}
We see from the more complete diagram
\[
\begin{tikzcd}[column sep=6.5em, row sep=4.5em]
& X \times I \arrow[r, "f'"] \arrow[dd, dotted, bend left=45, "\mathrm{pr}_1"{pos=.60}] \arrow[dr, phantom, "\substack{\text{cocart.}\\ \text{and cart.}}"] & Z(f) \arrow[dd, dotted, bend left=60, "p_f"] \\
X \arrow[ur, dotted, "\mathrm{inj}_1^X"] \arrow[r, hook, "i_1^X"] & X \amalg X \arrow[u, hook, "i_X"'] \arrow[r, "f \coprod \mathrm{id}_X"] \arrow[dr, phantom, "\substack{\text{cocart.}\\ \text{and cart.}}"] & Y \amalg X \arrow[u, hook, "j"'] \\
& X \arrow[u, hook, "i_0^X"'] \arrow[r, "f"] & Y, \arrow[u, hook]
\end{tikzcd}
\]
\printedpage{147}
% head: A. Grothendieck : Dérivateurs, Ch. XII, Édition des Universités de Montpellier II et Paris VII
% foot: Version 30 mars 2003
\[
i_f = f' \circ \mathrm{inj}_1^X = f' i_X i_1^X = j \ \ \underbrace{(f \amalg \mathrm{id}_X) i_1^X}_{\substack{\text{canonical inclusion}\\ \text{of } X \text{ into}\\ Y \coprod X}} \ \ ,
\]
it is \texttt{[an]} open Dwyer immersion, as a composite of $Y \longrightarrow Y \amalg X$ and $j : Y \amalg X \longrightarrow Z(f)$ which are such ($j$ obtained by cobase change from $i_X$, which is an open Dwyer immersion).

On the other hand, $X \xrightarrow{\mathrm{inj}_0^X} X \times I$ is \texttt{[an]} open Dwyer immersion and in $W$, hence couniversally in $W$, hence $Y \longrightarrow Z(f)$ is also in $W$, hence also the retraction $p_f$.

On the other hand, the proposition dual to
\[
\underline{\mathrm{Fib}}_W S \cap W_S = W^{\mathrm{u}}_{S\,\mathrm{univ}}
\]
is valid for $W$-cofibrations:

\msp{page 231}
\textbf{Proposition 16} (\footnote[125]{Has nothing to do with Cat - a general thing when there are enough $W$-cofibrations.}). \textit{On the full subcategory $\underline{\mathrm{Cofib}}_W(S)$ of $S\backslash\mathrm{Cat}$, the localizer ${}_SW$ is equal to $({}_SW)^{\mathrm{univ}}$ (couniversal $W$-equivalences), i.e. we have this: Let}
\[
S \longrightarrow X, \qquad S \longrightarrow Y
\]
\textit{be $W$-cofibrations, and let $f : X \longrightarrow Y$ be a morphism of $S\backslash\mathrm{Cat}$ such that $f \in W$. Then for every cobase change $S \longrightarrow S'$, the corresponding arrow}
\[
f' : X' = X \amalg_S S' \longrightarrow Y' = Y \amalg_S S'
\]
\textit{is in $W$.}

\textsc{Proof}. We factor $S \longrightarrow S'$ into
\[
S \xrightarrow{i} S'_0 \xrightarrow{j} S',
\]
with $i$ an open Dwyer immersion (hence a $W$-cofibration), and $j \in W$. So we have the diagram with cocartesian squares
\[
\begin{tikzcd}
S \arrow[r] \arrow[d, hook, "i"'] & X \arrow[r, "f"] \arrow[d, hook] & Y \arrow[d, hook] \\
S'_0 \arrow[r] \arrow[d, "\in W"'] & X'_0 \arrow[r, "f'_0", "W"'] \arrow[d, "W"] & Y'_0 \arrow[d, "W"] \\
S' \arrow[r] & X' \arrow[r, "f'"] & Y'.
\end{tikzcd}
\]
Since $i$ is a $W$-cofibration and $f \in W$, we have $f'_0 \in W$. On the other hand, since $S \longrightarrow X$, hence $S'_0 \longrightarrow X'_0$ is
\printedpage{148}
% head: A. Grothendieck : Dérivateurs, Ch. XII, Édition des Universités de Montpellier II et Paris VII
% foot: Version 30 mars 2003
\msp{page 232}
a $W$-cofibration, $X'_0 \longrightarrow X'$ is in $W$, and for the same reason, $Y'_0 \longrightarrow Y$ is too. So $f'$ is too. % [sic] Y'_0 \longrightarrow Y -> Y'_0 \longrightarrow Y', the diagram has Y'_0 \to Y'

\textbf{NB} Here, this is a formal argument, valid (in dual form) in any $(\mathcal{M}, W)$ where we have a factorization statement for an arbitrary arrow $f : X \longrightarrow Y$ into $X \xrightarrow{i_f} \bar{X} \xrightarrow{\alpha} Y$, with $i_f \in W$ and $\alpha$ a $W$-fibration. Then, if $X$, $Y$ are $W$-fibered over $S$, every $S$-arrow from $X$ to $Y$ that is in $W$, is in $W^{\mathrm{u}}_S$. And dually, if every arrow $X \longrightarrow Y$ factors as $X \longrightarrow \tilde{Y} \longrightarrow Y$, with $X \longrightarrow \tilde{Y}$ a $W$-cofibration and $\tilde{Y} \longrightarrow Y$ in $W$.

A (small) worry comes to me here: have I at least proved that if
\[
S \longrightarrow X
\]
is a $W$-cofibration, then for every arrow $S' \longrightarrow S$ in Cat, the cocartesian square in Cat % [sic] S' \longrightarrow S -> S \longrightarrow S', the square below has S \to S'
\[
\begin{tikzcd}
S \arrow[r] \arrow[d] & X \arrow[d] \\
S' \arrow[r] & X'
\end{tikzcd}
\]
is also $W$-cocartesian (\footnote[126]{in the old sense, special to Cat (def. 3, p. 131). To compare it with the general theory, cf. prop. 20, p. 262.})? Otherwise,
\msp{page 233}
what good would that square be?! To prove it, let us again factor $S \longrightarrow S'$ into
\[
\underset{\substack{\text{open}\\ \text{Dwyer}\\ \text{immersion}}}{S \hookrightarrow} S'_0 \underset{\in W}{\longrightarrow} S',
\]
whence the diagram of 2 cocartesian squares
\[
\begin{tikzcd}
S \arrow[r, "\mathrm{Cof}_W"] \arrow[d, hook, "D_W"'] & X \arrow[d, hook, "D_W"] \\
S'_0 \arrow[r, "\mathrm{Cof}_W"] \arrow[d, "\in W"'] & X'_0 \arrow[d, "\in W"] \\
S' \arrow[r, "\mathrm{Cof}_W"] & X',
\end{tikzcd}
\]
where $X'_0 \longrightarrow X'$ is $\in W$ by the hypothesis on $X$ under $S$. So we must verify:
\printedpage{149}
% head: A. Grothendieck : Dérivateurs, Ch. XII, Édition des Universités de Montpellier II et Paris VII
% foot: Version 30 mars 2003
\begin{enumerate}
\item[a)] A cocartesian square built on an arrow which is \texttt{[a]} Dwyer immersion, is $W$-cocartesian.
\item[b)] A commutative square in which two opposite sides are in $W$, is $W$-cocartesian.
\item[c)] The composite of two $W$-cocartesian squares is $W$-cocartesian.
\end{enumerate}

Let us provisionally admit b) and c), and prove a), for a cocartesian square
\[
\begin{tikzcd}
U \arrow[r, hook, "i"] \arrow[d, "p"'] & X \arrow[d, "q"] \\
U' \arrow[r, hook, "i'"] & X',
\end{tikzcd}
\]
where $i$ is an open Dwyer immersion. We proceed again as before, by factoring $p$ into $U \overset{j}{\hookrightarrow} U'_0 \xrightarrow{p'} U'$, where $j$ is \texttt{[an]} open Dwyer immersion, $p' \in W$. So we have a diagram of two adjoining cocartesian squares
\msp{page 234}
\[
\begin{tikzcd}
U \arrow[r, hook, "i \in D_W"] \arrow[d, hook, "D_W"'] & X \arrow[d, hook] \\
U'_0 \arrow[r, hook, "D_W"] \arrow[d, "p' \in W"'] & X'_0 \arrow[d, "q' \in W"] \\
U' \arrow[r, hook] & X',
\end{tikzcd}
\]
where the first square is made up of open immersions. We then know that it is $W$-cocartesian. In the second, we have $p' \in W$, hence $q' \in W$ since $i$ is a Dwyer immersion, hence a $W$-cofibration. So it is also $W$-cocartesian. We then conclude thanks to c) above. % [sic] immersions ouverts -> immersions ouvertes, agreement

\noindent Let us prove b), for
\begin{tikzcd}[baseline=(\tikzcdmatrixname-1-1.base)]
X_0 \arrow[r, "i_1"] \arrow[d, "i_2 \in W"'] & X_1, \arrow[d, "f_1 \in W"] \\
X_2 \arrow[r, "f_2"] & Y
\end{tikzcd}
by inserting it into
\begin{tikzcd}[baseline=(\tikzcdmatrixname-1-1.base)]
X_0 \arrow[r, "i_1"] \arrow[d, "i_2 \in W"'] & X_1, \arrow[d, "j_1"] \arrow[ddr, "f_1"] & \\
X_2 \arrow[r, "j_2"] \arrow[drr, "f_2"'] & \mathcal{X} \arrow[dr, "\varphi"] & \\
& & Y
\end{tikzcd}
where $\mathcal{X} = \displaystyle\int X_0$
\begin{tikzcd}[baseline=(\tikzcdmatrixname-2-1.base), row sep=small]
& X_1 \\
{} \arrow[ur] \arrow[dr] & \\
& X_2
\end{tikzcd}
.

Since $i_2 \in W$, we know that it follows that $j_1 \in W$, and since $f_1 \in W$, it follows that $\varphi \in W$, q.e.d.